\documentclass[a4paper,12pt]{article}

\usepackage[utf8]{inputenc}
\usepackage[german]{babel}
\usepackage{geometry}
\usepackage{hyperref}
\usepackage{graphicx}

\geometry{a4paper, margin=2.5cm}

\title{Projektkonzept: Hackintosh Companion App}
\author{Kilian Ketelhohn}
\date{\today}

\begin{document}

\maketitle

\section{Einleitung}
Im Rahmen des Moduls Mobile Programmierung wird eine Cross-Plattform Anwendung mit Flutter entwickelt. Ziel dieses Projekts ist die Erstellung einer mobilen Datenbank-Anwendung zur Verwaltung und Darstellung von Hackintosh-Konfigurationen.

Die Anwendung soll es ermöglichen, Informationen zu verschiedenen Nicht-Apple-Geräten (z. B. ThinkPads) übersichtlich darzustellen und schnell auf relevante Konfigurationshinweise zuzugreifen.

\section{Projektidee}
Die Hackintosh Companion App ist eine mobile Anwendung, die als Wissensdatenbank für die Installation von macOS auf nicht unterstützter Hardware dient.

Funktionen:
\begin{itemize}
\item Verwaltung von Geräten (z. B. ThinkPad Modelle)
\item Anzeige benötigter Kext-Treiber
\item Speicherung von BIOS- und OpenCore-Konfigurationen
\item Dokumentation von Problemen und Lösungen
\item Visuelle Darstellung von Hardware-Komponenten
\end{itemize}

\section{Technologien}

\subsection{Grundtechnologien}
\begin{itemize}
\item Flutter (Cross-Plattform Framework)
\item Dart (Programmiersprache)
\item SQLite (lokale Datenbank)
\item Android Studio (Entwicklungsumgebung)
\end{itemize}

\subsection{Verwendete Libraries}
\begin{itemize}
\item sqflite (SQLite Integration)
\item fl\_chart (Diagramme)
\item image\_picker (Kamera)
\item flutter\_3d\_controller (3D Darstellung)
\item http (REST API Kommunikation)
\end{itemize}

\section{REST API Integration}

\subsection{Motivation}
Neben der lokalen SQLite-Datenbank wird eine REST API integriert, um Daten zentral bereitzustellen und zwischen mehreren Nutzern synchronisieren zu können.

Dies ermöglicht eine skalierbare Architektur sowie zukünftige Erweiterungen wie Community-Funktionen und Cloud-Synchronisation.

\subsection{Grundlagen}
REST (Representational State Transfer) ist ein Architekturstil für Webservices, der über HTTP arbeitet und meist JSON als Datenformat verwendet.

HTTP-Methoden:
\begin{itemize}
\item GET: Daten abrufen
\item POST: Daten erstellen
\item PUT: Daten aktualisieren
\item DELETE: Daten löschen
\end{itemize}

\subsection{Architektur}
Die Anwendung folgt einem Client-Server-Modell:
\begin{itemize}
\item Client: Flutter App
\item Server: REST API Backend
\end{itemize}

\subsection{API Endpunkte}

\textbf{Devices:}
\begin{itemize}
\item GET /devices
\item GET /devices/{id}
\item POST /devices
\item PUT /devices/{id}
\item DELETE /devices/{id}
\end{itemize}

\textbf{Kexts:}
\begin{itemize}
\item GET /kexts
\item POST /kexts
\end{itemize}

\textbf{Issues:}
\begin{itemize}
\item GET /issues
\item POST /issues
\end{itemize}

\subsection{JSON Beispiel}
\begin{verbatim}
{
  "id": 1,
  "name": "ThinkPad T480",
  "cpu": "Intel i5",
  "gpu": "Intel UHD 620",
  "compatible": true
}
\end{verbatim}

\subsection{Flutter Integration}
\begin{verbatim}
final response = await http.get(
  Uri.parse('https://api.example.com/devices')
);
\end{verbatim}

\subsection{Offline-First Ansatz}
Die Anwendung kombiniert:
\begin{itemize}
\item lokale SQLite-Datenbank (Offline-Nutzung)
\item REST API (Online-Synchronisation)
\end{itemize}

\subsection{Synchronisation}
\begin{itemize}
\item Daten werden lokal gespeichert
\item Synchronisation bei bestehender Internetverbindung
\item Konfliktlösung über Zeitstempel
\end{itemize}

\subsection{Sicherheit}
\begin{itemize}
\item HTTPS
\item Optional: JWT Authentifizierung
\end{itemize}

\section{Datenbankdesign}

\subsection{devices}
\begin{itemize}
\item id
\item name
\item cpu
\item gpu
\item macOS kompatibel
\end{itemize}

\subsection{kexts}
\begin{itemize}
\item id
\item name
\item version
\item beschreibung
\end{itemize}

\subsection{device\_kexts}
\begin{itemize}
\item device\_id
\item kext\_id
\item notwendig
\item optional
\end{itemize}

\subsection{issues}
\begin{itemize}
\item device\_id
\item problem
\item loesung
\end{itemize}

\subsection{configs}
\begin{itemize}
\item device\_id
\item bios\_settings
\item opencore\_version
\item guide\_link
\end{itemize}

\section{App-Struktur}

\subsection{Home Screen}
\begin{itemize}
\item Geräteübersicht (ListView)
\item Suchfunktion
\end{itemize}

\subsection{Detail Screen}
\begin{itemize}
\item Geräteinformationen
\item Kexts
\item Probleme und Lösungen
\end{itemize}

\subsection{Visual Screen}
\begin{itemize}
\item 3D Modell eines Geräts
\item Interaktive Hardware-Komponenten
\end{itemize}

\subsection{Vergleichsansicht}
Vergleich von zwei Geräten hinsichtlich Kompatibilität und Konfiguration.

\subsection{Export}
Export von Daten als CSV oder XML.

\section{Animationen und Computergrafik}

\subsection{Boot-Animation}
Darstellung des Hackintosh Bootprozesses:
\begin{itemize}
\item BIOS
\item OpenCore Bootloader
\item macOS Start
\end{itemize}

\subsection{Hardware-Visualisierung}
\begin{itemize}
\item Interaktive Darstellung von Komponenten
\item Visualisierung von Problemen
\end{itemize}

\subsection{Diagramme}
\begin{itemize}
\item Erfolgsraten
\item Stabilitätsanalysen
\end{itemize}

\section{Sensorintegration}
\begin{itemize}
\item Kamera zur Erfassung von Konfigurationen
\item QR-Code Scanner
\end{itemize}

\section{Erweiterungsmöglichkeiten}
\begin{itemize}
\item Favoriten
\item Dark Mode
\item Community Features
\item Cloud Synchronisation
\end{itemize}

\section{Projektziel}
Ziel ist die Entwicklung einer funktionalen, benutzerfreundlichen und visuell ansprechenden Anwendung mit Fokus auf:

\begin{itemize}
\item Datenbankintegration
\item REST Kommunikation
\item Benutzeroberfläche
\item Animation und Visualisierung
\end{itemize}

\section{Fazit}
Die Anwendung verbindet moderne App-Entwicklung mit Datenbanktechnologien und Computergrafik und stellt eine praxisnahe Lösung für Hackintosh-Nutzer dar.

\end{document}